\begin{problem}{TST 2025 - Bài 2}{standard input}{standard output}{1 second}{1024 megabytes}

Bạn được cho một số nguyên dương $N$. Xét dãy số $A$ gồm $N$ phần tử thỏa mãn $A_i=1$ với mọi $i$ từ $1$ đến $N$. Nhiệm vụ của bạn là thực hiện một số thao tác sao cho sau khi thực hiện, với mọi $i$ từ $1$ đến $N$ phải thỏa mãn $A_i=i^i$. Bạn được phép thực hiện hai loại thao tác sau:

\begin{itemize}
    \item \texttt{+ i j k} ($1\le i,j,k\le N$): Cập nhật $A_k := A_i + A_j$.
    \item \texttt{* i j k} ($1\le i,j,k\le N$): Cập nhật $A_k := A_i \times A_j$.
\end{itemize}

\textbf{Lưu ý:} Trong một thao tác, các giá trị $i,j,k$ không nhất thiết phải đôi một phân biệt mà có thể bằng nhau.

\InputFile
Gồm duy nhất một số nguyên dương $N\ (1\le N\le 2\times 10^4)$.


\OutputFile
Dòng đầu tiên gồm một số nguyên không âm $M$ là số thao tác thực hiện $(0\le M\le 4\times 10^5)$.

Sau đó là $M$ dòng, mỗi dòng mô tả một trong hai thao tác ở trên.

\Scoring
\begin{center}
  \begin{tabular}{|c|c|l|} \hline
    \textbf{Subtask} & \textbf{Điểm} & \textbf{Giới hạn} \\ \hline
    1 & $5$ & $N\le 5$ \\ \hline
    2 & $15$ & $N\le 50$ \\ \hline
    3 & $80$ & Không có ràng buộc gì thêm \\ \hline
  \end{tabular}
\end{center}

\textbf{Cách tính điểm}

Giả sử test có tối đa $1$ điểm. Gọi $Q=\frac{M}N$. Khi đó, điểm của bạn sẽ được tính theo công thức như sau:
\begin{itemize}
   \item Nếu $Q>20$ thì bạn được $0$ điểm.
   \item Nếu $5<Q\le 20$ thì bạn được $0.9^{Q-5}$ điểm.
   \item Nếu $Q\le 5$ thì bạn được $1$ điểm.
\end{itemize}

Điểm của mỗi subtask là điểm thấp nhất của một test trong subtask đó.

\Example

\begin{example}
\exmpfile{example.01}{example.01.a}%
\end{example}

\end{problem}

